\newcommand{\titulus}{\nomenFesti{S. Bernardi, Abbatis \& Ecclesiæ Doctoris.}
\dies{Die 20. Augusti.}}
\newcommand{\oratio}{\pars{Oratio.}

\noindent Deus, qui beátum Bernárdum, abbátem, zelo domus tuæ succénsum, in Ecclésia tua lucére simul et ardére fecísti, eius nobis intercessióne concéde, ut, eódem spíritu fervéntes, tamquam fílii lucis iúgiter ambulémus.

\pars{Pro pace in universo mundo.} \scriptura{Sir. 50, 25; 2 Esdr. 4, 20; \textbf{H416}}

\vspace{-4mm}

\antiphona{II D}{temporalia/ant-dapacemdomine.gtex}

\vfill

\noindent Deus, a quo sancta desidéria, recta consília et iusta sunt ópera: da servis tuis illam, quam mundus dare non potest, pacem; ut et corda nostra mandátis tuis dédita, et hóstium subláta formídine, témpora sint tua protectióne tranquílla.

\noindent Per Dóminum nostrum Iesum Christum, Fílium tuum, qui tecum vivit et regnat in unitáte Spíritus Sancti, Deus, per ómnia sǽcula sæculórum.

\noindent \Rbardot{} Amen.}
\newcommand{\invitatorium}{\pars{Invitatorium.}

\vspace{-4mm}

\antiphona{IV}{temporalia/inv-fontemsapientiae.gtex}}
\newcommand{\hymnusmatutinum}{\pars{Hymnus.}

\vspace{-5mm}

\antiphona{V}{temporalia/hym-BernhardusDoctor.gtex}}
\newcommand{\matversus}{\noindent \Vbardot{} Non cessámus pro vobis orántes et postulántes.

\noindent \Rbardot{} Ut impleámini agnitióne voluntátis Dei.}
\newcommand{\lectioi}{\pars{Lectio I.} \scriptura{Is. 11, 1-9}

\noindent De libro Isaíæ prophétæ.

\noindent Hæc dicit Dóminus Deus:

\noindent Egrediétur virga de stirpe Iesse, et flos de radíce eius ascéndet;

\noindent et requiéscet super eum spíritus Dómini: spíritus sapiéntiæ et intelléctus, spíritus consílii et fortitúdinis, spíritus sciéntiæ et timóris Dómini; et delíciæ eius in timóre Dómini.

\noindent Non secúndum visiónem oculórum iudicábit neque secúndum audítum áurium decérnet;

\noindent sed iudicábit in iustítia páuperes et decérnet in æquitáte pro mansuétis terræ;

\noindent et percútiet terram virga oris sui et spíritu labiórum suórum interfíciet ímpium.

\noindent Et erit iustítia cíngulum lumbórum eius, et fides cinctórium renum eius.

\noindent Habitábit lupus cum agno, et pardus cum hædo accubábit;

\noindent vítulus et leo simul saginabúntur, et puer párvulus minábit eos.

\noindent Vítula et ursus pascéntur, simul accubábunt cátuli eórum; et leo sicut bos cómedet páleas.

\noindent Et ludet infans ab úbere super forámine áspidis; et in cavérnam réguli, qui ablactátus fúerit, manum suam mittet.

\noindent Non nocébunt et non occídent in univérso monte sancto meo, quia plena erit terra sciéntia Dómini, sicut aquæ mare opériunt.}
\newcommand{\responsoriumi}{\pars{Responsorium 1.} \scriptura{\Rbar{} Is. 1, 11 \Vbar{} Is. 45, 8; \textbf{H36}}

\vspace{-5mm}

\responsorium{I}{temporalia/resp-egredieturvirga-CROCHU.gtex}{}}
\newcommand{\lectioii}{\pars{Lectio II.} \scriptura{Is. 11, 10-16}

\noindent In die illa radix Iesse stat in signum populórum; ipsam gentes requírent, et erit sedes eius gloriósa.

\noindent Et erit in die illa: rursus exténdet Dóminus manum suam ad possidéndum resíduum pópuli sui, quod relíctum erit ab Assýria et ab Ægýpto et a Phátros et ab Æthiópia et ab Elam et a Sénnaar et ab Emath et ab ínsulis maris;

\noindent et levábit signum in natiónes et congregábit prófugos Israel et dispérsos Iudæ cólliget a quáttuor plagis terræ.

\noindent Et auferétur zelus Ephraim, et hostes Iudæ abscindéntur; Ephraim non æmulábitur Iudam, et Iudas non pugnábit contra Ephraim.

\noindent Et volábunt in úmeros Philísthim ad mare, simul prædabúntur fílios oriéntis; in Edom et Moab exténdent manus suas, et fílii Ammon obœ́dient eis.

\noindent Et exsiccábit Dóminus linguam maris Ægýpti et levábit manum suam super flumen in fortitúdine spíritus sui et percútiet illud in septem rivos, ita ut transíre fáciat eos calceátos.

\noindent Et erit via resíduo pópulo meo, qui relinquétur ab Assýria, sicut fuit Israéli in die illa, qua ascéndit de terra Ægýpti».}
\newcommand{\responsoriumii}{\pars{Responsorium 2.} \scriptura{\Rbar{} Is. 11, 10 \Vbar{} Cf. Hab. 3, 3; \textbf{H29}}

\vspace{-5mm}

\responsorium{VIII}{temporalia/resp-ecceradixjesse-CROCHU.gtex}{}}
\newcommand{\lectioiii}{\pars{Lectio III.} \scriptura{Sermo 83, 4-6: Opera omnia, Edit. Cisterc. 2 [1958], 300-302}

\noindent Ex Sermónibus sancti Bernárdi abbátis super Cánticum canticórum.

\noindent Amor per se súfficit, is per se placet, et propter se. Ipse méritum, ipse prǽmium est sibi. Amor præter se non requírit causam, non fructum: fructus eius, usus eius. Amo quia amo; amo ut amem. Magna res amor, si tamen ad suum recúrrat princípium, si suæ orígini rédditus, si refúsus suo fonti, semper ex eo sumat unde iúgiter fluat. Solus est amor ex ómnibus ánimæ mótibus, sénsibus atque afféctibus, in quo potest creatúra, etsi non ex æquo, respondére Auctóri, vel de símili mútuam repéndere vicem. Nam cum amat Deus, non áliud vult, quam amári: quippe non ad áliud amat, nisi ut amétur, sciens ipso amóre beátos, qui se amáverint.

\noindent Sponsi amor, immo Sponsus amor solam amóris vicem requírit et fidem. Líceat proínde redamáre diléctam. Quidni amet sponsa, et sponsa Amóris? Quidni amétur Amor?

\noindent Mérito cunctis renúntians affectiónibus áliis, soli et tota incúmbit amóri, quæ ipsi respondére amóri habet in reddéndo amórem. Nam et cum se totam effúderit in amórem, quantum est hoc ad illíus fontis perénne proflúvium? Non plane pari ubertáte fluunt amans et Amor, ánima et Verbum, sponsa et Sponsus, Creátor et creatúra, non magis quam sítiens et fons.

\noindent Quid ergo? Períbit propter hoc, et ex toto evacuábitur nuptúræ votum, desidérium suspirántis, amántis ardor, præsuméntis fidúcia, quia non valet ex æquo cúrrere cum gigánte, dulcédine cum melle conténdere, lenitáte cum agno, candóre cum lílio claritáte cum sole, caritáte cum eo qui cáritas est? Non. Nam etsi minus díligit creatúra, quóniam minor est, tamen si ex tota se díligit, nihil deest ubi totum est. Proptérea sic amáre, nupsísse est, quóniam non potest sic dilígere, et parum dilécta esse, ut in consénsu duórum íntegrum stet perfectúmque connúbium. Nisi quis dúbitet, ánimam a Verbo et prius amári, et plus.}
\newcommand{\responsoriumiii}{\pars{Responsorium 3.} \scriptura{\Rbardot{} Sir. 45, 9 \Vbardot{} 1 Th. 5, 8; \textbf{H378}}

\vspace{-5mm}

\responsorium{II}{temporalia/resp-amaviteumdominus-CROCHU-cumdox.gtex}{}}
\newcommand{\hymnuslaudes}{\pars{Hymnus.}

\vspace{-5mm}

\antiphona{V}{temporalia/hym-ArcanaSacrae.gtex}}
\newcommand{\lectiobrevis}{\pars{Lectio Brevis.} \scriptura{Sap. 7, 13-14}

\noindent Sapiéntiam sine fictióne dídici et sine invídia commúnico; divítias illíus non abscóndo. Infinítus enim thesáurus est homínibus; quem qui acquisiérunt, ad amicítiam in Deum se paravérunt propter disciplínæ dona commendáti.}
\newcommand{\responsoriumbreve}{\pars{Responsorium breve.} \scriptura{Sir. 45, 9}

\antiphona{VI}{temporalia/resp-amaviteum.gtex}}
\newcommand{\preces}{\noindent Christum Deum sanctum, fratres, exaltémus, orántes ut serviámus illi in sanctitáte et iustítia coram ipso ómnibus diébus nostris,~\gredagger{} et acclamémus:

\Rbardot{} Tu solus sanctus, Dómine.

\noindent Qui tentári voluísti per ómnia pro similitúdine nostra absque peccáto,~\gredagger{} miserére nostri, Dómine Iesu.

\Rbardot{} Tu solus sanctus, Dómine.

\noindent Qui nos omnes ad perfectiónem caritátis vocásti,~\gredagger{} sanctífica nos, Dómine Iesu.

\Rbardot{} Tu solus sanctus, Dómine.

\noindent Qui nos iussísti esse salem terræ et lucem mundi,~\gredagger{} illúmina nos, Dómine Iesu.

\Rbardot{} Tu solus sanctus, Dómine.

\noindent Qui voluísti ministráre, non ministrári,~\gredagger{} fac nos tibi et frátribus humíliter servíre, Dómine Iesu.

\Rbardot{} Tu solus sanctus, Dómine.

\noindent Tu, splendor glóriæ Patris et figúra substántiæ eius,~\gredagger{} fac ut in glória vultum tuum respiciámus, Dómine Iesu.

\Rbardot{} Tu solus sanctus, Dómine.}
\newcommand{\benedictus}{\pars{Canticum Zachariæ.} \scriptura{Cf. Eccli. 39, 8}

\vspace{-4mm}

\antiphona{VII a}{temporalia/ant-replevitsanctumsuum.gtex}

\vspace{-2mm}

\scriptura{Lc. 1, 68-79}

\vspace{-2mm}

\cantusSineNeumas
\initiumpsalmi{temporalia/benedictus-initium-vii-a-auto.gtex}

%\vspace{-1.5mm}

\input{temporalia/benedictus-vii-a.tex} \Abardot{}}
\newcommand{\benedicamuslaudes}{\cuminitiali{}{temporalia/benedicamus-memoria-laudes.gtex}}
\include{hebdomadaxx}
\include{feriav}
